\section{Geometric systems and classification by bases}
\label{sys:section-geometric-systems}

Throughout this section the distinguished model $M_0$, its fixed pointed
pattern, and the pregeometry $P$ are those of
Proposition~\ref{geom:pointed-levels} and
Theorem~\ref{geom:pregeometry}.
All model sizes exceed the threshold of
Corollary~\ref{geom:tail-no-growth}.  Thus a strong inclusion of pointed
models with equal $P$-sets is equality.  At an eligible uniform level
$\mu$, all vertices have size $\mu$.  A statement about unions of
objects at that level concerns chains whose total size is $\mu$, as in
the working conditions of Definition~\ref{lev:working-package}.

We first prove structural and presentation assertions without assuming
categoricity of systems.  The classification argument then proceeds by
the number of axes.  Its first assertion extends a root isomorphism and
prescribed bijections of the other basis pieces; this gives categoricity
for the current shape.  That categoricity supplies the working package
used to prove the second assertion: prescribed bijections of the basis
pieces added to a fixed system extend to an isomorphism over that system.
By splitting off the last axis, the second assertion gives the first at
the next dimension.  We establish the conditional package before carrying
out this paired induction.

\subsection{Shapes, relative bases, and arrows}

\begin{definition}[Admissible shapes and geometric systems]
\label{sys:shapes}
Let $A$ be a finite linearly ordered set of axes.  An admissible family
$L\subseteq\mathcal P(A)$ of \emph{active labels} is defined recursively.
On no axes the family is $\{\varnothing\}$.  If $k$ is a new last axis,
choose an upward closed subset $Q$ of the previously constructed family
$L$, and put
\[
 L'=L\cup\{s\cup\{k\}:s\in Q\}.
\]
Here upward closed means that $s\in Q$, $s\subseteq t$, and $t\in L$
imply $t\in Q$.

A geometric system of shape $L$ is a family
$(X_t:t\subseteq A)$ of pointed models in a common top model, with
$X_t\leq_\K X_u$ when $t\subseteq u$, satisfying the following condition.
There is a partition $(H_s:s\in L)$ of a basis of $P(X_A)$ such that
$H_s\subseteq P(X_s)$ and
\begin{equation}
\label{sys:vertex-span}
 P(X_t)=\cl_{P(X_A)}\left(\bigcup_{s\in L,\ s\subseteq t}H_s\right)
 \qquad(t\subseteq A).
\end{equation}
At uniform level $\mu$ we require $|H_s|=\mu$ for every $s\in L$.
The resulting scheme is denoted $\mathcal G_L$.
The basis partition witnesses the definition; it is not named in the
system structure.

For a vertex $t$, write
\[
 D_t(X)=\cl_{P(X_A)}\left(\bigcup_{r\subsetneq t}P(X_r)\right).
\]
For $t=\varnothing$ this is the closure of the empty set in the
pregeometry.  An arrow $X\to Y$ is a commuting family of pointed strong
embeddings for which a witnessing partition for $X$ extends to a
witnessing partition for $Y$.  Equivalently, after identifying $X$ with
its image, each finite tuple from $P(X_t)$ independent over $D_t(X)$
remains independent over $D_t(Y)$.
\end{definition}

The equivalence in the last sentence, including its independence from
the initially chosen partition, will be proved below.  All closures and
independence assertions are computed in a common containing pointed
model.  Theorem~\ref{geom:pregeometry} makes this convention
independent of the chosen common strong extension.
The notation $\cl_{P(X_A)}$ means the closure $\cl_{X_A}$ defined
there; its subscript emphasizes the pregeometry in which it is
computed.  We omit that subscript when the ambient model is clear.

\begin{lemma}[Combinatorics of shapes]
\label{sys:shape-combinatorics}
An admissible family contains $\varnothing$ and is closed under unions.
Its restriction to any subset of the axes is admissible, with the
inherited order.  Consequently every $t\subseteq A$ has a largest
active subset
\[
 a_L(t)=\bigcup\{s\in L:s\subseteq t\}.
\]
In a geometric system, $X_t=X_{a_L(t)}$ whenever $t$ is inactive.
\end{lemma}

\begin{proof}
The union assertion follows by induction on the axes.  Unions of old
labels are old.  If a union involves $s\cup\{k\}$ with $s\in Q$, its
old part is an active label containing $s$, and therefore belongs to
$Q$.  This also handles the union of two labels containing $k$.
Restriction to a set of axes replaces $L$ and $Q$, at every retained
step, by their restrictions.  The restricted $Q$ is upward closed in
the restricted $L$.  Steps on discarded axes make no contribution.
This proves the restriction assertion.

The displayed union is an active label, since the family is finite and
contains $\varnothing$.  The collections of active labels below $t$
and below $a_L(t)$ agree.  Equation~\eqref{sys:vertex-span} therefore
gives equality of their $P$-sets.  The inclusion
$X_{a_L(t)}\leq_\K X_t$ and
Corollary~\ref{geom:tail-no-growth} give equality of the models.
\end{proof}

\begin{remark}
Union closure alone would not imply the recursive last-axis
description.  For example, $\{\varnothing,\{1,2\}\}$ is union closed,
but its two layers on splitting off axis $2$ have different active
families.  We use the recursive definition, including the upward
closed set $Q$, when identifying a system with an arrow between
systems on fewer axes.
\end{remark}

\begin{lemma}[Replacement and extension of relative bases]
\label{sys:relative-bases}
Suppose $X$ is a geometric system.  For each active $s$, choose any
basis $B_s$ of $P(X_s)$ over $D_s(X)$.  Then the $B_s$ form a
witnessing partition for $X$.

A componentwise strong system embedding is an arrow precisely when
it preserves the relative independence in
Definition~\ref{sys:shapes}.  For an arrow, every witnessing partition
of its domain extends to one of its codomain.
\end{lemma}

\begin{proof}
Fix a witnessing partition $(H_s:s\in L)$, and enumerate $L$ in an
order extending inclusion.  For active $s$,
\[
 D_s(X)=\cl\left(\bigcup_{u\in L,\ u\subsetneq s}H_u\right).
\]
When the piece at $s$ is replaced, previously processed labels not
below $s$ are independent, together with the original piece at $s$,
over this lower span: they are parts of a single independent set.
Replacing $H_s$ by any basis of the same flat over its lower span
preserves both independence and the span of the union.  Explicitly,
let $C$ be the current pieces below $s$.  Then
$\cl(C\cup B_s)=\cl(C\cup H_s)=P(X_s)$, and $C\cup B_s$ is
independent.  The remaining current pieces were independent over
$\cl(C\cup H_s)$, since the full current partition was independent.
They remain independent over the equal flat $\cl(C\cup B_s)$.
Thus the new full partition is independent and has the same span.
Applying this replacement at the finitely many active labels gives the
first assertion and Equation~\eqref{sys:vertex-span} for the new
partition.

If pieces extend along an embedding, any finite independence
calculation in a vertex over its lower span is the same calculation
using the corresponding pieces in the larger independent set.
Relative independence is therefore preserved.

Conversely, let $(B_s)$ be an arbitrary witnessing partition of the
domain, viewed as a subsystem of $Y$.  Preservation implies that
$B_s$ is independent over $D_s(Y)$.  Extend it to a relative basis
$B_s\cup C_s$ of $P(Y_s)$ over $D_s(Y)$.  The first assertion, applied
to $Y$, says that these choices assemble into a witnessing partition
of $Y$.  This proves both the equivalence and the final assertion.
\end{proof}

\begin{lemma}[Structural properties and the last-axis description]
\label{sys:structural}
Arrows compose and satisfy coherence.  Same-level chains of total
size the level have unions in $\mathcal G_L$, with the common-upper-
bound clause also for arrows into a larger eligible level.

Restriction to any subset of axes preserves systems and arrows.
If $L'=L\cup\{s\cup\{k\}:s\in Q\}$, splitting off the last axis
identifies a system of shape $L'$ with an arrow $X\to Y$ in
$\mathcal G_L$ such that the complementary piece at $s$ has size
$\mu$ for $s\in Q$ and is empty for $s\notin Q$.
\end{lemma}

\begin{proof}
Composition follows from the relative independence criterion.
For coherence, suppose $X\subseteq Y$, and both embed by arrows in
$Z$.  Independence over $D_t(X)$ implies independence over $D_t(Z)$,
and hence over its subset $D_t(Y)$.  Componentwise coherence is the
AEC coherence axiom.

For a continuous chain, extend witnessing pieces successively using
Lemma~\ref{sys:relative-bases}, and take their unions at limit stages.
Every finite independent sublist and every finite dependence witness
occurs at a stage.  The resulting union pieces are independent and
span every vertex of the componentwise union.  Each contains $\mu$
elements from the first stage and has size at most $\mu$, so all
pieces still have size $\mu$.  For a chain which is not initially
continuous, insert its earlier unions: the preceding construction
applies recursively, and the common-upper-bound argument in the next
paragraph supplies its connecting arrows.

If each stage maps by an arrow into $Z$, a finite tuple independent
over the lower span of the union is independent over the lower span
at any stage containing it.  Its image is independent over $D_t(Z)$.
The componentwise AEC common-upper-bound axiom supplies the strong
maps, so the union map is an arrow.  This argument also permits a
larger level for $Z$.

On restricting to axes $B\subseteq A$, retain the pieces whose labels
are contained in $B$.  They are a basis of $P(X_B)$ and satisfy the
required span identities.  For a retained vertex all its proper
subsets are still retained; its relative independence test is
therefore exactly the old test.

Finally let $Z$ have shape $L'$, and put $X_s=Z_s$ and
$Y_s=Z_{s\cup\{k\}}$.  A partition for $Y$ is obtained by combining
$H_s$ with $H_{s\cup\{k\}}$ when $s\in Q$, and using $H_s$ alone
otherwise.  These choices extend the pieces of $X$, and their
complements have the asserted sizes.  Conversely, extend pieces of
$X$ to pieces of $Y$ along such an arrow, give the old pieces labels
$s$ and their complements labels $s\cup\{k\}$, and use $X$ and $Y$
as the two layers.  The span identities give a system of shape $L'$.
Inactive vertices cause no exception, by
Lemma~\ref{sys:shape-combinatorics}.
\end{proof}

For later use, preservation of relative independence also preserves
finite rank calculations modulo the lower span.  For finite tuples
$a,c$ in a vertex, the condition that $a$ be independent over
$D_t(X)\cup c$ is determined by the ranks, modulo $D_t(X)$, of
sublists of $a,c$.  These ranks are determined by which sublists are
independent.  Thus arrows preserve this condition as well.

\subsection{Order presentations and their rank}

\begin{proposition}[Tiled presentations]
\label{sys:tiling}
Let $L$ be admissible and let $\mu$ be sufficiently large for the
pointed pregeometry and the no-proper-pair conclusion.  The order
template $E$ gives a presentation in $\mathcal G_L$ as follows.
Choose a fixed initial order carrying the whole enumeration of $M_0$.
For each active label $t$, choose a tile $T_t$, preceded by a buffer
$B_t$, and concatenate these finitely many pairs after the initial
order, in any fixed order.  Every tile and buffer is a dense linear
order without endpoints of size $\mu$.  In vertex $s$ retain the
initial order, all buffers, and just the tiles with labels $t\subseteq
s$.

These presentations have the following properties.
\begin{enumerate}
\item Each active relative piece has rank $\mu$.
\item Order embeddings of the corresponding tiles and buffers,
fixing the initial pointed data, induce arrows, including between
different eligible levels.
\item A presentation has an extension by an arrow gaining $\mu$
independent complementary coordinates at every active label
simultaneously.
\item Tiles may be chosen with $\mu$ pairwise disjoint convex
substrips, each of size $\mu$.  Further insertions of such substrips
preserve the preceding assertions and the naturalness of the
underlying $\K$-diagrams.
\end{enumerate}
No categoricity assumption about $\mathcal G_L$ is used here.
\end{proposition}

\begin{proof}
Let $I_s$ be the indicated suborder and set $X_s=E(I_s)$, using the
literal embeddings supplied by Theorem~\ref{found:order-template}.
All $I_s$ have size $\mu$, since they contain all buffers.  The
models have size $\mu$ and are strongly included according to their
labels.  The initial order was chosen to contain the supports of the
fixed enumeration of $M_0$.  Its unchanged term pattern carries the
same pointed model in every vertex, and coherence gives the required
strong inclusion of $M_0$.

Choose a relative basis $H_s$ of $P(X_s)$ over the union of its
proper lower $P$-sets for every active $s$.  We first prove that the
indexed family of these choices is independent; this also proves
that the pieces are disjoint.  Suppose there is a finite circuit
of labeled entries, allowing a two-entry equality dependence when
two differently labeled entries coincide.  Choose a maximal label $s$ represented among its
entries.  Retain the finite term supports of all entries with label
$s$, and indeed retain every support position in a tile available
at $s$.  In each unavailable tile replace all support positions by
distinct positions in its preceding buffer, beyond all the finitely
many retained positions in that buffer.  Such positions exist because
the buffer has no greatest point and is dense.  No retained position
lies in that unavailable tile.  Consequently this replacement
preserves the order and equality pattern of all circuit supports and
fixes the whole initial pointed order.

The labels and this order pattern preserve the pointed finite
geometry comparisons.  One can formulate this either using the
natural $\K$-diagrams or by taking the term model on the fixed initial
order and the finite supports and applying its two pointed strong
embeddings into the original presentation.  In the latter formulation
Theorem~\ref{geom:pregeometry} preserves and reflects the
independence relation.

An entry with label $t\ne s$ in the circuit has no label properly
containing $s$, by maximality.  Thus $s\nsubseteq t$.  After the
replacement its term uses only common buffers, initial data, and
tiles with labels contained in $s\cap t$.  Its image belongs to
$P(X_{s\cap t})\subseteq D_s(X)$.  The entries with label $s$ have
not moved.  The preserved circuit makes a sublist of $H_s$ dependent
over $D_s(X)$, a contradiction.  The union is therefore independent.
It spans every vertex by induction on the labels: its relative basis
spans that vertex over lower vertices, which have already been
spanned.  If a vertex is inactive, the order retained there is exactly
the order retained at its largest active subset.  This proves the
decomposition without first assuming the system axioms.

We give the fresh-skeleton observation needed for the rank assertion.
If $J\subseteq I$ and $i\in I\setminus J$, the skeleton entry at $i$
does not belong to $E(J)$.  Otherwise it equals a labeled term on a
finite tuple of positions from $J$.  Enlarge the order by a second
position $i'$ having the same order pattern as $i$ relative to that
finite tuple.  Literal pattern preservation would make both
distinguished skeleton entries equal to that term.  This contradicts
their required distinctness.  The enlargement need only duplicate
the cut relative to the finite tuple, not the cut relative to all of
$J$.

For an active $t$, omit $T_t$ from $I_t$.  The resulting order still
has size $\mu$, contains the pointed data and all proper lower
vertices, and presents a strong submodel.  Enumerate $\mu$ points of
$T_t$ and adjoin them one at a time.  The fresh-skeleton observation
makes each successor inclusion proper.  All these models have size
$\mu$, because the buffers were retained.  By
Corollary~\ref{geom:tail-no-growth}, each successor adds a point of
$P$.  Choose one such point at each successor.  Relative closedness
of $P$ in strong submodels makes these choices independent over the
initial omitted-tile model.  There are $\mu$ choices, so the rank over
the proper lower vertices is at least $\mu$.  It is at most $\mu$
by the size of $X_t$.  This proves the first assertion, including the
root case $t=\varnothing$.

Next consider componentwise order embeddings into enlarged tiles and
buffers, identifying the source with its image.  Suppose a finite
tuple $a$ in $P(X_s)$, independent over $D_s(X)$, becomes dependent
over $D_s(Y)$.  Dependence has a finite witness consisting of points
from proper lower vertices of $Y$.  Fix the finitely many old
positions supporting $a$.  Move every other position needed for the
witness into its corresponding old tile or buffer, preserving its
order relative to the fixed positions.  Density and absence of
endpoints of the old order suffice for each finite pattern; it is
not necessary that the old order be dense in the enlarged order.
The initial pointed order remains fixed.  Each witness remains in
its designated proper lower vertex, now in $X$, and its finite
dependence diagram is preserved.  This contradicts independence of
$a$ over $D_s(X)$.  Lemma~\ref{sys:relative-bases} therefore gives an
arrow.  The proof did not require equal source and target sizes.

To obtain simultaneous gains, enlarge all tiles with $\mu$ fresh
points and, if desired, enlarge the buffers.  For a particular active
$t$, begin the strict-growth argument over the order containing the
old $t$-tile, all enlarged proper-lower tiles, all enlarged buffers,
and the initial data.  This base contains $X_t$ and all proper lower
vertices of $Y$.  Adjoining the fresh points of the $t$-tile yields
$\mu$ independent points over that larger base.  Hence the
complementary rank at $t$ is $\mu$.  The preceding arrow assertion
allows any old pieces to extend, and
Lemma~\ref{sys:relative-bases} combines the relative choices into
one partition.  This proves the simultaneous assertion.

Finally a lexicographic sum, indexed by a linear order of size
$\mu$, of dense endpoint-free orders of size $\mu$ is still dense
and endpoint-free and has size $\mu$.  Its summands are the desired
convex substrips.  In particular, for a presentation indexed by a
well-order $H$, one can use $H\times\Q$ and reserve intervals of
$H$ of order type and cardinality $\mu$.  Taking $H$ of order type
$\mu\cdot\mu$ supplies $\mu$ disjoint such intervals.  Additional
summands or strips do not change any of the preceding arguments.
Underlying $\K$-diagrams remain natural because the entire system
was obtained by designated term subcollections of the single
template $E$.
\end{proof}

\subsection{Obtaining the working package from categoricity}

\begin{proposition}[Conditional working package]
\label{sys:working-package}
Fix an admissible shape $L$.  Suppose that $\mathcal G_L$ is
categorical at every level in an unbounded class of eligible levels
closed sufficiently for the preceding constructions.  Assume also
the eligible-subclass hypothesis of Convention~\ref{lev:eligible-subclass}.
On an unbounded further class of sufficiently closed levels, it satisfies the full working
package of Definition~\ref{lev:working-package}, with its actual
system isomorphisms and its arrows from Definition~\ref{sys:shapes}.
The resulting diagrams, saturation, and calculus include mixed
arrows between eligible levels, and the output class retains the
eligible-subclass hypothesis.

The closure requirements are uniform in the finite shape and the
fixed pointed template.  They use only the smaller and larger input
levels specified by the general working-package construction.
\end{proposition}

\begin{proof}
The structural conditions are Lemma~\ref{sys:structural}.
Proposition~\ref{sys:tiling} supplies term presentations and mixed
presentation arrows before categoricity.  By the assumed
categoricity, every object at an input level is isomorphic to its
presentation.

We check the extra relative-basis condition throughout the reversal
argument.  In each tile and buffer use the order $H(\delta)$ from
Theorem~\ref{found:order-orbits}.  For a tuple of length $\theta$, record
its term labels, the fixed initial positions it uses, and, separately
in each of the finitely many variable orders, the order and
realized-cut record of its support positions.  If
$\theta<\cf(\delta)$, equal records are carried to one another by
order automorphisms of the respective tiles and buffers.  Their
product fixes the initial data and induces an automorphism of the
whole system.  The set of possible records has a bound depending
only on $\theta$, the finite shape, and the template.  Hence the
system orbits at that arity have a uniform set bound.

Number these orbits and record the arrows between them, together
with the partial map from support records to orbit names.  This is
again a bounded palette.  If reversal fails unboundedly often,
choose repeated profiles at eligible levels $a<\mu$ exactly as in
Theorem~\ref{found:reversal}, with $a^\theta\leq\mu$ and the required
large cofinalities for the orders.  Apply the ordinal-position
shrinking of Proposition~\ref{found:alignment} simultaneously to the
finitely many tiles and buffers.  Given at most $a$ parameters in
the large system, the resulting small presentation contains them,
has size $a$ in every tile and buffer, and computes the support
records of all its $\theta$-tuples correctly.  Its inclusion is an
actual arrow by Proposition~\ref{sys:tiling}.  Thus it is an aligned
small geometric system, not merely a componentwise subsystem.

For completeness, the contradiction uses two orbit names $p,r$ with
$p\to r$ and $r\not\to p$.  Repeatedly choose aligned small copies,
each containing the previous one, and process their at most
$a^\theta$ tuples in the large level.  Whenever a tuple can still be
sent to $p$, send it onward to $r$.  From that point it cannot reach
$p$.  Coherence makes the successive small copies a chain of
arrows.  After $\theta^+$ rounds every $\theta$-tuple of their union
occurred at a proper round.  Both unions exist by
Lemma~\ref{sys:structural}: the small union has total size $a$, and
the large construction has total size $\mu$.  Alignment says that
none of the tuples of the small union can reach $p$.  Categoricity
of that union at $a$ says that its orbit $p$ is realized, a
contradiction.  This proves eventual reversal at each fixed arity.

Mixed L\"owenheim--Skolem closure also has the required form.
In a presentation, close the supports of a prescribed small set
under finitely many intermediate and exterior positions in every
tile and buffer, repeat countably many times, and pad each order
to the chosen smaller eligible size.  The resulting dense
endpoint-free suborders give a small presentation containing the
set.  The initial pointed order is retained in full.  Its inclusion
is an arrow by Proposition~\ref{sys:tiling}.  Transport through a
presentation of an arbitrary large object gives mixed containment;
coherence supplies comparisons between such subobjects.

We now apply the first route of
Proposition~\ref{lev:scheme-diagram-routes}.  Its structural
hypothesis is Lemma~\ref{sys:structural}; its presentation and mixed
downward-closure hypotheses are Proposition~\ref{sys:tiling} and
the preceding paragraph.  Categoricity is the input assumption,
and the required arity-by-arity reversal and uniform orbit bounds
were proved above.  Thus the proposition applies to the actual
system arrows, including their relative-basis condition.

In particular, the canonical comparison at a smaller input level
is computed in a small geometric subobject.  Two such sections
are contained in a common small section by mixed downward closure,
and coherence gives its arrows, so the comparison is well defined.
The projections between comparisons at two input levels are onto
because the smaller system is a tiled presentation and extends by
the mixed order arrows of Proposition~\ref{sys:tiling}.  The bounded
palette therefore stabilizes, exactly as in the first route of
Proposition~\ref{lev:scheme-diagram-routes}.

Choose receiving levels which remain eligible, are strictly closed
under the resulting arity bounds, and have sufficiently large
smaller input levels below them for every bounded request.
These are the closure refinements described in
Section~\ref{lev:refining-levels}, supplied simultaneously in
Theorem~\ref{close:common-levels}; closure under the next input
level is not being used to infer eligibility by itself.  At such a
level, countable alternating matching in smaller input-level
sections gives actual system isomorphisms on equal bounded
diagrams, and realizes possible bounded joint diagrams over their
given subtuples.  Thus diagrams, homogeneity, and saturation have
the system interpretation required here, rather than merely their
componentwise $\K$ interpretation.  Order embeddings and order
amalgamation make the diagrams natural on term tuples, so
Proposition~\ref{lev:diagram-coding} also applies.

It remains to check the additional exhaustion condition, which
Proposition~\ref{lev:scheme-diagram-routes} explicitly leaves as a
separate obligation.  A finite tuple independent modulo the
proper lower vertices stays so when its target diagram agrees with
a tuple placed along an actual presentation arrow: the comparison
is realized by a system isomorphism, and such an isomorphism
preserves relative independence.  Every tuple in an exhaustive
sequence is handled at a stage containing it.  The finite
relative-independence criterion in Lemma~\ref{sys:relative-bases}
then supplies the additional arrow condition.  Componentwise
strongness is the conclusion of the same proposition's exhaustion
argument.  This completes the working-package verification.

All refinements used here are the bounds for the original reversal,
alignment, diagram stabilization, and saturation arguments, applied
to finitely many orders and the fixed initial data.  None introduces
an amalgamation assumption on the original AEC; the amalgamation
used below is the conclusion of the working package and
Theorem~\ref{calc:scheme-amalgamation}.
\end{proof}

\subsection{Fresh coordinates and the two induction assertions}

Fix pieces $(H_s^X:s\in L)$ in a system $X$.
For an arrow $X\to Y$, extend them to pieces
\[
 H_s^Y=H_s^X\mathbin{\dot\cup}C_s^Y.
\]
The sets $C_s^Y$ are its \emph{complementary pieces}.  Their sizes
may be any cardinals at most $\mu$, including zero, even though all
the total pieces have size $\mu$.

\begin{definition}[Joint freshness]
\label{sys:fresh}
A list of coordinates in an extension $Y$ of $X$, with individual
active labels, is jointly fresh over $X$ if it is a sublist of the
complementary pieces for some such extension of the pieces of $X$.
At a single label $s$, the corresponding condition is independence
over
\[
 \cl\bigl(P(X_s)\cup D_s(Y)\bigr).
\]
For several labels these relative conditions assemble into joint
freshness by Lemma~\ref{sys:relative-bases}.  Coordinates are
regarded as occurrences in their designated vertex sorts.
\end{definition}

Freshness is preserved by arrows of extensions.  Indeed every
failure or success of finite independence over the displayed flat
is determined by finite rank calculations using parameters from
$P(X_s)$ and the lower span.  The observation following
Lemma~\ref{sys:structural} preserves those calculations.  This also
shows that freshness is an invariant of the type over $X$ once the
working package is available.

\begin{theorem}[Classification with prescribed basis maps]
\label{sys:basis-classification}
For each finite admissible shape $L$, at all sufficiently closed
uniform levels $\mu$, the following assertions hold.
\begin{enumerate}
\item If $X,Y\in\mathcal G_L$, an isomorphism
$e_\varnothing:X_\varnothing\simeq Y_\varnothing$ and arbitrary
bijections between the chosen nonroot pieces of $X$ and $Y$ extend
to a system isomorphism $X\simeq Y$.
In particular $\mathcal G_L$ is categorical at these levels.
\item Suppose $X\to Y$ and $X\to Z$ are arrows in
$\mathcal G_L$.  Choose pieces of $X$, extend them to pieces of
$Y$ and $Z$, and prescribe a bijection
$C_s^Y\simeq C_s^Z$ for every active label $s$.
Then there is a system isomorphism $Y\simeq_X Z$ realizing every
prescribed bijection.  The complementary cardinalities may be
less than $\mu$, and zero is allowed.
\end{enumerate}
The level requirements are uniform over the finitely many shapes
at any fixed number of axes.
\end{theorem}

\begin{proof}
We prove the two assertions in order, by the number of axes.
On no axes, the first assertion is exactly the assertion that a
prescribed isomorphism of the root models is an isomorphism of the
systems.  Existence of such a root isomorphism, hence categoricity
of the zero-axis scheme, is
Proposition~\ref{geom:pointed-levels}.

Suppose the second assertion is known for all shapes on fewer
axes.  For a shape with last-axis description $L'$, apply the first
assertion for $L$ to its lower layers and then the second assertion
for $L$ to its upper layers.  The pieces added at labels in $Q$
are exactly the pieces bearing labels $s\cup\{k\}$ in $L'$.
Prescribe their given bijections, and prescribe the empty maps at
the other labels.  This produces the desired isomorphism of the
two-layer systems by Lemma~\ref{sys:structural}.
If an auxiliary root basis has been chosen, replace the target
root basis by its image under the prescribed root isomorphism;
Lemma~\ref{sys:relative-bases} permits that replacement without
changing any prescribed nonroot piece.  Thus the first assertion
at the next number of axes follows from the preceding second
assertion.

It remains to prove the second assertion for a fixed shape $L$
whose first assertion has now been established.  Refine the
levels using Proposition~\ref{sys:working-package}.  All the
one-dimensional calculus and dominated-hull results are then
available for $\mathcal G_L$.  The rest of the proof is given in
four claims.

\begin{claim}[Uniqueness of finite fresh types]
\label{sys:finite-fresh-types}
For a fixed base $X$ and a finite ordered list of active labels,
all jointly fresh tuples with those labels have the same type
over $X$.
\end{claim}

\begin{proof}
At zero axes this is the base-change calculation of
Proposition~\ref{geom:base-change}.  Types in the pointed scheme
over a whole pointed base are the underlying $\K$-types: an
amalgam or type identification over that base already fixes the
point.  All objects at the level have full $P$-rank by
Corollary~\ref{geom:tail-no-growth}.

At a positive number of axes suppose two finite jointly fresh
tuples $b^0,b^1$ have different types over $X$.  Bounded testing
gives a bounded tuple $f$ from $X$ such that
\[
 D(f,b^0)\ne D(f,b^1).
\]
Choose a cardinal $\chi<\mu$ large enough for the orbit bound at
arity $|f|$, for the finite tuple lengths, and for the template
labels, with $\chi^{|f|+\aleph_0}=\chi$.  Let $\epsilon$ be the
least cardinal for which $2^\epsilon>\chi$.  Then
$\epsilon\leq\chi$, there are at most $\chi$ nodes at all proper
levels of the binary tree of height $\epsilon$, and
$2^\epsilon\leq2^\chi<\mu$.

Build a tree of extensions of $X$ at level $\mu$.  At a node $V$,
the arrow from $X$ preserves the diagram of $f$.  Bounded-tuple
homogeneity therefore supplies an isomorphism $X\simeq V$
fixing $f$.  Transport the two witness extensions along this
isomorphism to form the children of $V$.  Their marked tuples
are fresh over $V$, and their diagrams with $f$ have the two
specified different values.  At limit nodes take unions along
the branch.  Extend chosen basis pieces at every step.  Hence
along every branch all marked coordinates, at their respective
labels, lie in complementary pieces over the initial $X$.
The total size along a branch is $\mu$, so all unions are
allowed.  Notice that the return isomorphism fixing $f$ uses
diagram homogeneity, not categoricity alone.

At every leaf make a further rich extension by
Proposition~\ref{sys:tiling}, so that the root has complementary
rank $\mu$ over the original $X_\varnothing$.  There is a common
system $U$ with the same property: present any system and make a
rich extension, transporting it over $X$ using categoricity.
For each level $\eta<\epsilon$ reserve in $U$ a finite tuple
$c_\eta$ of coordinates with the same labels as the marked
tuple at that level.  Choose all of these coordinates from
complementary pieces over $X$.  Their total number is at most
$\chi$.

For a leaf $W$, its root marked coordinates and those in $U$
are independent over $P(X_\varnothing)$.  Both full root
complements have size $\mu$, and the prescribed match uses
fewer than $\mu$ coordinates.  Extend that match to a bijection
of root complements.  The already proved second assertion for
zero axes gives
\[
 W_\varnothing\simeq_{X_\varnothing}U_\varnothing
\]
carrying all root marked coordinates correctly.  At nonroot
labels extend the prescribed matches of marked coordinates to
bijections of the full pieces, whose sizes are $\mu$.  The
first assertion for the present shape extends these maps and
the root isomorphism to an isomorphism $W\simeq U$.
The old nonroot pieces of $X$ need not be fixed in this step.

Let $f_W$ be the image of $f$ in $U$.  Two different branches
first separate at a level $\eta$; preservation of diagrams
along their arrows gives different values of
$D(f_W,c_\eta)$ after the two leaf isomorphisms.  Thus the
$2^\epsilon$ images of $f$ lie in different orbits over the
fixed set of at most $\chi$ reserved coordinates.  The
parameter-orbit bound of Lemma~\ref{calc:parameter-orbits},
applied to the working scheme, is $\chi$ at
these arities, a contradiction.
\end{proof}

For each active $s$, let $p_s^X$ denote the unique fresh singleton
type given by the claim.  Such types exist by the rich extensions
of Proposition~\ref{sys:tiling} and categoricity.  Their realizations
lie outside the base vertex, so they are nonalgebraic.

\begin{claim}[Free lifts preserve freshness]
\label{sys:free-fresh}
The free lift of $p_s^X$ along an arrow $X\to V$ is fresh over $V$.
Finite jointly fresh tuples have the product of their singleton
types.  Bounded jointly fresh lists have a unique whole type
specified by their labels and admit all-dominated hulls.
\end{claim}

\begin{proof}
Let $a$ realize the free lift in an extension $Y$ of $V$.
If $a$ is not fresh, finite character gives a finite tuple
$c\subseteq P(V_s)$ with
\[
 a\in\cl\bigl(c\cup D_s(Y)\bigr).
\]
This condition is invariant under system isomorphisms and is
therefore determined by the bounded diagram of $(a,c)$ in $Y$.
Its witnesses in lower vertices need not lie in $V$; the
isomorphism-invariant condition itself is what the diagram
records.

Put $m=|c|+1$, and choose an $m$-tuple $b$ jointly fresh over
$X$, all at label $s$.  Every singleton marginal is $p_s^X$ by
Claim~\ref{sys:finite-fresh-types}.  Take a sufficiently long
bounded family of approximations, inside $X$, for
$\tp(b/X)$, as in Lemma~\ref{calc:approximations}.
By the support-copy construction in that Lemma, every
approximating tuple has the same diagram as
$b$ without parameters, and is consequently independent
modulo $D_s(X)$.  In the free lift to $V$, each coordinate has
the same marginal as $a$.  For each coordinate, the majority
estimate gives the diagram of $(a,c)$ outside a uniformly
bounded set of approximation indices.  Choose the family
longer than the union of these finitely many exceptional sets.
One approximating tuple $(a_1,\ldots,a_m)$ is independent
modulo $D_s(X)$ and has, coordinate by coordinate, the required
diagram with $c$ in $V$.

Since $X\to V$ is an arrow, this tuple remains independent
modulo $D_s(V)$.  The comparison just obtained is an equality of
$\mathcal G_L$-diagrams,
\[
 D_{\mathcal G_L}(a,c;Y)=D_{\mathcal G_L}(a_i,c;V).
\]
The established working package therefore gives a system isomorphism
$F_i:Y\cong V$ taking $a$ to $a_i$ and fixing the displayed tuple $c$.
It carries every lower vertex onto its counterpart and preserves the
pointed closure, so $F_i[D_s(Y)]=D_s(V)$.  Transporting the dependence
witness gives $a_i\in\cl(c\cup D_s(V))$ for every $i$.  The latter
flat has relative rank at most $|c|$, whereas the tuple has
relative rank $m>|c|$.  This contradiction proves that the
free lift is fresh.

Realize a finite product successively by free lifts along a
chain of extensions.  The assertion just proved puts each
new coordinate in a complementary piece over the preceding
system.  Extending pieces along the chain makes the whole
tuple jointly fresh over $X$.  Claim~\ref{sys:finite-fresh-types}
then says that every jointly fresh finite tuple has this
product type.  A bounded jointly fresh list is consequently
product-independent in the sense of
Theorem~\ref{iso:independent-lists}.  That theorem gives
its unique whole type and its witnessing free chains; its
all-dominated-hull conclusion gives existence over the list.
\end{proof}

We now know the types of fresh coordinate lists and have all-dominated
hulls over them.  To extend a partially prescribed basis map, we also
need to cover an arbitrary bounded tuple using only boundedly many
new coordinates.  The next claim provides this size control, including
when the target basis has already been chosen.

\begin{claim}[Bounded rank cost and chosen coordinates]
\label{sys:rank-cost}
Let $X\to Y$ be an arrow.  A finite tuple $d$ of $Y$ is
contained in a subextension $X\to U\to Y$ whose complementary
pieces have finite sizes.  For an arbitrary bounded tuple
$d$, their sizes can be bounded by
$\max(|d|,\aleph_0)$.

If specified pieces of $Y$ extend those of $X$, the subextension
containing a bounded $d$ can be chosen with all its
complementary pieces subsets of the specified complementary
pieces of $Y$, still of bounded total size.
\end{claim}

\begin{proof}
For finite $d$, take an all-dominated object over $X+d$ and
embed it into $Y$ by Proposition~\ref{iso:prime}.  Call its
image $U$.  Suppose one complementary piece of $U$ is
infinite.  Choose a sequence $(a_i:i<\omega)$ from that piece.
It is jointly fresh and hence product-independent over $X$.
The witnessing-chain assertion of
Theorem~\ref{iso:independent-lists} retains this
configuration in a chain
\[
 X=V_0\longrightarrow V_1\longrightarrow\cdots
 \longrightarrow V_\omega,
\]
where $a_i$ is free over $V_i$ from $X$ and belongs to
$V_{i+1}$.  We may construct this chain inside the actual $U$.  At each stage,
finite domination makes the actual next coordinate $a_i$ free over
$V_i$ from $X$, since its products with the preceding coordinates
remain those of the original list.  Choose an all-dominated increment
over $V_i+a_i$ and use primeness to embed it into $U$, fixing this
base and this coordinate.  At limits take unions inside $U$, using
the common-upper-bound clause.  Thus $V_\omega\leq U$, and the tuple
$d$ and every $a_i$ retain their actual positions in $U$.  Equivalently,
the witnessing-chain amalgamation may retain all of these data in a
common extension.

The tuple $d$ is finite, so Proposition~\ref{calc:local-character}
at the limit of
cofinality $\omega$ makes the actual type $\tp(d/V_\omega;U)$ free from
some $V_n$.  Hence $d$ and $a_n$ have product over $V_n$.
The base-change rule of Proposition~\ref{calc:products},
together with the freeness
of $a_n$ over $V_n$ from $X$, gives that $a_n$ and $d$ have
product over $X$.  Domination of $U$ over $X+d$ now makes
$\tp(a_n/U)$ free from $X$.  This is impossible by
Claim~\ref{sys:free-fresh}: the free lift of its fresh
singleton marginal must be fresh over $U$, whereas
$a_n\in P(U_s)$.  Therefore every complementary piece of
$U$ is finite.

For a bounded list of cardinality $\tau$, include its entries
successively by this finite construction, with the preceding
subextension as base, and take unions at limit stages.
There are finitely many labels.  At most $\tau$ finite
increments therefore add at most
$\kappa=\max(\tau,\aleph_0)<\mu$ coordinates in total.
All objects continue to have size $\mu$, and their unions have
total size $\mu$.  This proves the bounded version.

For the final assertion fix the target partition
$H_s^X\mathbin{\dot\cup}C_s^Y$.  Starting with a bounded-rank
subextension containing $d$, repeat the following construction
countably many times.  Choose relative bases for the current
subextension over $X$.  Each gained basis element in vertex
$t$ has a finite support in the fixed target basis, consisting
of old coordinates and coordinates of $C_s^Y$ with
$s\subseteq t$.  Request all the latter supporting
coordinates in their own vertices, and include them in the
next subextension using the bounded construction just proved.
At every step at most $\kappa$ coordinates are requested and
the total new rank remains at most $\kappa$.

In the union let $C_s$ consist of the requested coordinates
of label $s$.  These sets are independent, together with the
pieces of $X$, because they are subsets of the target
partition.  They span every union vertex over its old pieces:
an element occurs at a stage, is spanned there by the gained
pieces and the old pieces, and the finite target supports of
those gained pieces were requested at the next stage.
Conversely the closed $P$-set of each union vertex contains
all the indicated coordinates below its label.  Thus the
span identities hold in both directions.  The $C_s$ are
precisely complementary pieces of the union, have total
size at most $\kappa$, and are subsets of the prescribed
$C_s^Y$.
\end{proof}

\begin{claim}[Primeness over bounded complementary pieces]
\label{sys:bounded-prime}
If $X\to U$ has a bounded list $b$ enumerating all its
complementary pieces, then $U$ is all-dominated over $X+b$.
It embeds over $X$ into any extension realizing a prescribed
matching of these coordinates.
\end{claim}

\begin{proof}
The list $b$ is product-independent by
Claim~\ref{sys:free-fresh}.  Take an all-dominated object over
$X+b$ and embed it into $U$ over these data.  Its image $V$
contains in each vertex the old $P$-set and all the
complementary coordinates with labels below that vertex.
Relative closedness and the span identity imply
$P(V_t)=P(U_t)$ for every $t$.  The no-proper-pair conclusion
then gives $V_t=U_t$ for every vertex.  Thus $V=U$, proving
the domination assertion.

A prescribed matching of $b$ into another extension gives
the same whole type over $X$, by
Claim~\ref{sys:free-fresh}.  Primeness of the all-dominated
object therefore gives the desired arrow fixing $X$ and
realizing that matching.
\end{proof}

We finish the second assertion of the theorem.  Identify the two
bases with $X$, and keep the prescribed bijections of all
complementary pieces.  Choose countable exhaustions of the two
targets by bounded lists; these exist because the eligible level
$\mu$ has cofinality $\omega$ and there are only finitely many
vertex sorts.  Construct alternating isomorphisms between
subextensions of the two targets.  Initially use $\id_X$.
Maintain the invariant that the complementary pieces in each
current subextension are subsets of the specified target pieces,
and that the current isomorphism agrees with their prescribed
bijections.

At a forth step, apply Claim~\ref{sys:rank-cost} over the current
subextension on the first side to include the next bounded list
and to use only remaining specified coordinates.  Its new
complementary list is bounded.  The corresponding coordinates
on the second side are jointly fresh over its current
subextension.  Claim~\ref{sys:bounded-prime}, transported by the
current isomorphism, embeds the enlarged subextension into the
second target over that isomorphism and realizes their
prescribed matching.  Its image is the next subextension there.
The span identities show that its complementary pieces are
exactly the specified images of the new pieces, so the invariant
is preserved.  A back step applies the same construction with
the roles interchanged.

The union of these countably many coherent isomorphisms covers
both targets, since every list in both exhaustions is eventually
included.  It is therefore a system isomorphism over $X$.
It realizes all prescribed basis bijections: every basis element
is eventually in the domain, and the invariant determines its
image.  This proves the second assertion, with no lower bound
on complementary cardinalities.

The induction is complete.  At each fixed finite number of axes
there are only finitely many active-label families.  Taking the
maximum of their finitely many bounds gives the stated
uniformity.  Simultaneous closure through all finite numbers of
axes is addressed separately in
Theorem~\ref{close:common-levels}.
\end{proof}
